\section{Methodology}

Our sparse coupling hypothesis requires measuring all 10 dimension pairs (from 5 dimensions: truthfulness, robustness, fairness, safety, privacy) to distinguish broad coupling (many pairs significant) from narrow coupling (few pairs). Difficulty-independence requires explicit control via partial correlation. This section explains why each design choice solves the coupling characterization problem.

Returning to our medical diagnosis scenario---independent evaluation reports 92\% truthfulness, 89\% fairness, but phi coefficient on instance-level co-occurrence reveals whether the 8\% truthfulness failures overlap with the 11\% fairness failures. If phi = 0.45, approximately 20\% of failures are compound failures invisible to per-dimension scores. This coupling detection is the methodology's core capability.

\subsection{Phi Coefficient for Coupling Strength}

We measure coupling via the phi coefficient ($\phi$), an effect size for $2\times2$ contingency tables. For each dimension pair (A, B) and model, we construct a contingency table from instance-level binary labels (pass/fail):

\begin{table}[h]
\centering
\begin{tabular}{l|cc}
 & B\_pass & B\_fail \\
\hline
A\_pass & $n_{11}$ & $n_{12}$ \\
A\_fail & $n_{21}$ & $n_{22}$ \\
\end{tabular}
\end{table}

Phi coefficient quantifies association strength:

\[
\phi = \frac{n_{11} \cdot n_{22} - n_{12} \cdot n_{21}}{\sqrt{(n_{11}+n_{12})(n_{21}+n_{22})(n_{11}+n_{21})(n_{12}+n_{22})}}
\]

Phi ranges from 0 (independence) to 1 (perfect coupling), providing an interpretable scale: $\phi \geq 0.3$ indicates medium effect size in social science conventions. We pair this with chi-square significance tests ($p < 0.01$ threshold) to separate real coupling from sampling noise.

\textbf{Why phi over alternatives?} Odds ratios are less interpretable (exponential scale), chi-square alone provides no effect size, and Pearson correlation assumes continuous data. Phi directly measures co-occurrence probability in binary outcomes---exactly our coupling definition. For the medical scenario, phi quantifies how much knowing a patient failed truthfulness increases probability of fairness failure.

\subsection{Partial Correlation for Difficulty Control}

Instance difficulty confounds coupling estimates. Hard instances may fail on all dimensions simultaneously, creating spurious correlation independent of shared vulnerabilities. We control this via partial correlation, computing phi after partialing out difficulty:

\[
\phi_{\text{partial}}(A, B \mid \text{difficulty}) = \text{correlation residual after regressing } A, B \text{ on difficulty}
\]

We operationalize difficulty as a composite score derived from model confidence (simulated in Phase 4; API logprobs in production). Figure~\ref{fig:difficulty_independence} validates the control strategy: difficulty shows $|r| < 0.2$ with all dimensions, confirming statistical independence required for partial correlation.

\textbf{Why partial correlation over stratification?} We use both as dual validation. Partial correlation provides continuous control across the difficulty spectrum, while quartile stratification (Figure~\ref{fig:quartile_stratified}) offers robustness checking---coupling must persist in $\geq$3 quartiles. The dual approach guards against distributional assumptions in either method. For the medical scenario, this ensures the coupling we detect isn't simply ``hard minority-group cases fail on everything.''

\subsection{Mantel Test for Model-Specific Fingerprints}

Model-specific coupling profiles require comparing coupling matrices across model pairs. Each model produces a $5\times5$ symmetric coupling matrix (10 unique pairs). The Mantel test measures matrix similarity via permutation-based correlation:

\begin{enumerate}
\item Compute Pearson $r$ between vectorized matrices (GPT-4 vs Claude-3)
\item Permute rows/columns of one matrix 10,000 times
\item Calculate $r$ for each permutation
\item $p$-value = proportion of permutations with $|r| \geq$ observed $|r|$
\end{enumerate}

\textbf{Why Mantel over element-wise correlation?} Mantel preserves matrix structure---dimensions have inherent ordering (truthfulness, robustness, fairness, safety, privacy). Element-wise correlation treats pairs as independent, ignoring positional information. Our threshold $r < 0.7$ defines ``distinct profiles''---Bonferroni-corrected $p < 0.0167$ confirms significance.

\subsection{Bonferroni Correction for Multi-Pair Analysis}

Testing $\geq$3 significant pairs per model (h-m2) requires multiple comparison correction. We use Bonferroni correction for family-wise error rate control:

\[
\alpha_{\text{adjusted}} = 0.01 / 30 \approx 0.00033 \quad \text{(10 pairs $\times$ 3 models)}
\]

Each pair must meet both $\phi \geq 0.3$ AND $p < \alpha_{\text{adjusted}}$ to count toward the $\geq$3 threshold. This conservative approach controls false positives but may exclude borderline pairs.

\textbf{Why Bonferroni over FDR?} Bonferroni guarantees family-wise error rate $< 0.01$---no false positives across all 30 tests. FDR (Benjamini-Hochberg) controls expected false discovery proportion, accepting some false positives for higher power. We chose Bonferroni for confirmatory analysis (h-m2 tests a specific prediction: $\geq$3 pairs), reserving FDR for exploratory sensitivity analysis.

\subsection{Intuition Building via Figures}

Our methodology produces three key visualizations:

\textbf{Figure~\ref{fig:difficulty_independence} (difficulty\_independence.png):} Correlation heatmap showing difficulty vs. each dimension. All $|r| < 0.2$ validates that difficulty is orthogonal to dimension-specific failures---required assumption for partial correlation.

\textbf{Figure~\ref{fig:partial_vs_raw} (partial\_vs\_raw\_phi.png):} Scatter plot comparing raw phi (x-axis) vs partial phi (y-axis) for 6 significant pairs across 3 models. Points above the diagonal (5/6 cases) show retention $>100\%$---difficulty acts as suppressor, not confounder. This surprising result indicates difficulty masks true coupling strength rather than inflating it.

\textbf{Figure~\ref{fig:quartile_stratified} (quartile\_stratified\_phi.png):} Heatmap showing coupling strength across difficulty quartiles (Q0=easy to Q3=hard) for truthfulness-robustness and fairness-safety pairs. Persistence in 3--4 quartiles per model validates coupling independence from difficulty gradients---complementary to partial correlation.

\subsection{Technical Depth Balance}

Main paper presents phi formula, partial correlation intuition, and Mantel test overview with visual aids. Appendix provides full contingency tables (30 tables: 10 pairs $\times$ 3 models), power analysis explaining why n=100/dimension underpowers h-c1 Mantel test (requires $n \geq 500$ for $r=0.5$ detection at 80\% power), and sensitivity analysis comparing Bonferroni vs FDR correction.

This methodology addresses three design requirements: (1) measure all 10 pairs to distinguish broad from sparse coupling, (2) control difficulty to isolate genuine coupling, (3) compare models to test fingerprint hypothesis. The next section reports synthetic validation results demonstrating methodology capabilities.
